\subsection{Semi-reflexive spaces}

DEFINITION 2. --- Let E be a locally convex space. We say that E is semi-reflexive if the canonical mapping \(c_{\mathbf{E}}\) from E into E'' is bijective.

This implies that E is Hausdorff, and that every linear form on \(E'\) , which is continuous for the strong topology \(\beta(E', E)\) , is of the form \(x' \mapsto \langle x, x' \rangle\) with \(x \in E\) , i.e.~continuous for the weak topology \(\sigma(E', E)\) .

THEOREM 1. --- A locally convex Hausdorff space is semi-reflexive if and only if every bounded subset of E is relatively compact for the weakened topology \(\sigma(E, E')\) . If E is semi-reflexive, the strong dual \(E_b'\) of E is barrelled.

The second assertion follows from prop. 3 (IV, p.~15), and the identity between bounded subsets for the initial topology and for the weakened topology of E (III, p.~27, cor. 3).

To say that E is semi-reflexive means that the topology on \(E_{b}^{\prime}\) is compatible with the duality between E and \(E^{\prime}\) , in other words, by Mackey's theorem (IV, p.~2, th. 1) that the topology on \(E_{b}^{\prime}\) is coarser than \(\tau(\mathrm{E}^{\prime}, \mathrm{E})\) (and in fact is identical with it); by definition (IV, p.~2), this means that every closed, convex and bounded subset of E is compact for \(\sigma(\mathrm{E}, \mathrm{E}^{\prime})\) , and this is equivalent to saying that every bounded subset of E is relatively compact for \(\sigma(\mathrm{E},\mathrm{E}^{\prime})\) , because the closed convex envelope of a bounded subset of E is bounded (III, p.~3, prop. 1).

COROLLARY. --- Let E be a locally convex semi-reflexive space. Every closed vector subspace M of E is semi-reflexive; moreover, the strong topology on \(E^{\prime}/M^{\circ}\) (considered as the dual of M) is the quotient of the strong topology on \(E^{\prime}\) .

Let B be a bounded subset of M. Since B is bounded in E, and the weakened topology \(\sigma(\mathbf{M},\mathbf{M}^{\prime})\) on M is induced by \(\sigma(\mathbf{E},\mathbf{E}^{\prime})\) (IV, p.~10, prop. 11), the closure of B in M endowed with \(\sigma(\mathbf{M},\mathbf{M}^{\prime})\) is compact. Hence, by th. 1, M is semi-reflexive. The last assertion of the corollary follows from prop. 10 of IV, p.~9, applied to the set S of all closed, convex and bounded subsets of E.

Remarks. --- 1) Suppose E is semi-reflexive. Every subset of E which is convex, closed and bounded for the initial topology is compact for the topology \(\sigma(E, E')\) (IV, p.~1, prop. 1). * On the other hand, the unit sphere (with the equation \(\|x\|=1\) ) of an infinite dimensional hilbertian space E is closed and bounded for the initial topology, but is not closed for the weakened topology, even if E is semi-reflexive. *

\begin{enumerate}
\def\labelenumi{\arabic{enumi})}
\setcounter{enumi}{1}
\item
  By remark 3 of IV, p.~5, we can reformulate th. 1 as follows: the Hausdorff space E is semi-reflexive if and only if it is quasi-complete for its weakened topology. If it is semi-reflexive, then it is quasi-complete for its initial topology (IV, p.~5, Remark 2).
\item
  Under the hypotheses of the above corollary, the space E/M is not necessarily semi-reflexive (IV, p.~63, exerc. 10).
\end{enumerate}

